% Answers to Chapter 3, from lectures/L03/appendix/c_solutions.md; the self-test from
% lectures/L03/README.md.

\facitavsnitt{c3}{Kapitel 3}{Linjära ekvationer och olikheter}
\begin{facit}

\svar{1.1} \sv{a} $4R + 14$ \sv{b} $U^2 + 6U = U(U + 6)$ \sv{c} $3I + 2$

\svar{1.2} \sv{a} $(R + 5)(R - 5)$ \sv{b} $3U(U + 2)$ \sv{c} $(I - 4)^2$

\svar{2.1} \sv{a} $15 = R \cdot 0{,}3$ ger $R = 50\,\Omega$.
\sv{b} $50 \times 0{,}3 = 15\,\text{V}$, vilket stämmer.

\svar{2.2} \sv{a} $24 = (R_1 + 40) \times 0{,}4$ ger $R_1 + 40 = 60$, alltså $R_1 = 20\,\Omega$.
\sv{b} Samma svar, $R_1 = 20\,\Omega$.
\sv{c} $I = \frac{24}{20 + 40} = 0{,}4\,\text{A}$, vilket stämmer.

\svar{2.3} \sv{a} $3 = \frac{3}{R_1 + 3} \times 9$ ger $R_1 = 6\,\Omega$.
\sv{b} $\frac{3}{6 + 3} \times 9 = 3\,\text{V}$, vilket stämmer.

\svar{2.4} $R = \frac{\tau}{C} = 500\,\Omega$

\svar{2.5} \sv{a} $U_{\text{CE}} \cdot I_{\text{C}} \leq P_{\max}$, det vill säga
$U_{\text{CE}} \cdot 0{,}05 \leq 0{,}5$ \sv{b} $U_{\text{CE}} \leq 10\,\text{V}$

\svar{2.6} \sv{a} $10{,}4\,\text{V}$
\sv{b} $12 - 0{,}8I \geq 9$ ger $I \leq 3{,}75\,\text{A}$; tecknet vänder vid divisionen med
$-0{,}8$.
\sv{c} $12 - 0{,}8I = 13 - 1{,}0I$ ger $I = 5\,\text{A}$, där båda ger $8\,\text{V}$.
\sv{d} Det första batteriet: $5{,}6\,\text{V}$ mot $5{,}0\,\text{V}$.

\svar{2.7} \sv{a} $-0{,}5\,\text{V} \leq \delta \leq 0{,}5\,\text{V}$
\sv{b} $4{,}5\,\text{V} \leq u \leq 5{,}5\,\text{V}$

\svar{2.8} \sv{a} $5\,\text{A}$ \sv{b} $I = 2\,\text{A}$ och $I = 8\,\text{A}$
\sv{c} $2\,\text{A} < I < 8\,\text{A}$

\svar{3.1} \sv{a} $x = 5$ \sv{b} $x = 9$ \sv{c} $x = 5$ \sv{d} $x = 24$ \sv{e} $x = 4$
\sv{f} $x = -6$

\svar{3.2} \sv{a} $x = 4$ \sv{b} $x = 2$ \sv{c} $x = 5$
\sv{d} Alla $x$: leden är identiska, så ekvationen är en identitet.
\sv{e} Ingen lösning: $x$ försvinner och kvar blir $8 = 5$, som är falskt.

\svar{3.3} \sv{a} $x = 6$ \sv{b} $x = 8$ \sv{c} $x = 10$ \sv{d} $x = 3$

\svar{3.4} \sv{a} $x = 6$ \sv{b} $x = 15$ \sv{c} $R_1 = 1\,\text{k}\Omega$
\sv{d} $\frac{6}{0{,}2} = \frac{U}{0{,}5}$ ger $U = 15\,\text{V}$.

\svar{3.5} \sv{a} $I = \frac{U}{R}$ \sv{b} $R = \frac{P}{I^2}$ \sv{c} $R_2 = R_{\text{TOT}} - R_1$
\sv{d} $C = \frac{\tau}{R}$ \sv{e} $R_{\text{i}} = \frac{E - U}{I}$
\sv{f} $P_{\text{in}} = \frac{P_{\text{ut}}}{\eta}$

\svar{3.6} \sv{a} $x < 5$: öppen ring vid $5$, pil åt vänster.
\sv{b} $x \geq 4$: fylld ring vid $4$, pil åt höger.
\sv{c} $x < -4$: öppen ring vid $-4$, pil åt vänster.
\sv{d} $x \geq 3$: fylld ring vid $3$, pil åt höger.
\sv{e} $x < -3$: öppen ring vid $-3$, pil åt vänster.

\svar{3.7} \sv{a} $1 < x < 6$ \sv{b} $-1 \leq x \leq 3$
\sv{c} $5I \leq 0{,}25$ ger $I \leq 0{,}05\,\text{A} = 50\,\text{mA}$.
\sv{d} $4{,}0 \leq x \leq 6{,}6$

\svar{3.8} \sv{a} $x = 6$ eller $x = -6$ \sv{b} $x = 8$ eller $x = -2$
\sv{c} $x = 3$ eller $x = -4$ \sv{d} $x = -4$, en enda lösning.
\sv{e} Ingen lösning: ett absolutbelopp är aldrig negativt.

\svar{3.9} \sv{a} $-2 \leq x - 5 \leq 2$, alltså $3 \leq x \leq 7$. \sv{b} $1 < x < 5$
\sv{c} $\abs{u - 5} \leq 0{,}1$, alltså $4{,}9\,\text{V} \leq u \leq 5{,}1\,\text{V}$.
\sv{d} Nej: $\abs{1015 - 1000} = 15 > 10$.

\svar{3.10} \sv{a} $120\,\Omega$ \sv{b} $45\,{}^{\circ}\text{C}$
\sv{c} Högst $95\,{}^{\circ}\text{C}$, eftersom $\Delta T \leq 75\,{}^{\circ}\text{C}$.
\sv{d} $\Delta T = \frac{R - R_0}{R_0\,\alpha}$

\svar[Testa dig själv]{T} \sv{1} $x = 8$; båda led blir då $24$.
\sv{2} $I = 0{,}03\,\text{A} = 30\,\text{mA}$ \sv{3} $I \leq 6\,\text{A}$

\end{facit}
